\subsection{Robust policy search with bounds}\label{\refsec:RobustRoot}
The politico-economic equilibria below are identified from first order conditions rather than by maximising political objectives directly. This keeps the state space small, since it lets the closed-form savings shares stand in for the distribution of past savings, but a root problem lacks the safeguards of the maximisation it came from: a root may be a minimum, there may be several, and the optimum may sit at a corner where no root exists. This subsection states the machinery that restores them, for a generic first order condition; it is implemented in \texttt{python/gridsearch}, and the model enters only through the function evaluated.

Let $f(\tau)$ denote a first order condition, in our applications the derivative of a political objective $\mathcal{W}$ with respect to the tax rate, that is safe to evaluate on $[l,u]$ with $0\leq l < u < 1$, and let $z(\tau) = f\left([\tau]_{[l,u]}\right)$ denote its bounded evaluation, where $[\tau]_{[l,u]} \equiv \min(u, \max(l,\tau))$.

\smalltitle{A bounded-root reparameterization.}
Gradient-based solvers search all of $\mathbb{R}$, while $f$ is only meaningful on $[l,u]$. Instead of solving $f(\tau) = 0$ subject to bounds, we solve the unconstrained auxiliary problem
\begin{align}\label{\refeq:root}
    h(z,\tau) = z - |z| \underbrace{\left[a_0\min\left(\tau - l,0\right) + a_1\max\left(\tau - u, 0\right)\right]}_{\equiv g(\tau)} = 0,
\end{align}
where $a_0, a_1>0$ are technical parameters. This gives three regions:
\begin{align*}
    h\left(z,\tau\right) = \left\lbrace \begin{array}{ll} f(l)-|f(l)| a_0 (\tau - l), & \tau<l \\
    f(\tau), & \tau\in[l,u] \\
    f(u)-|f(u)|a_1(\tau-u), & \tau>u\end{array}\right., && \dfrac{\partial h}{\partial \tau} = \left\lbrace \begin{array}{ll} -|f(l)|a_0, & \tau<l \\
    \dfrac{\partial f}{\partial \tau}, & \tau\in[l,u] \\
    -|f(u)|a_1, & \tau>u\end{array}\right.
\end{align*}
$f$ is never evaluated outside $[l,u]$, yet $h$ is defined everywhere, affine with a weakly negative slope outside the interval, so an unconstrained solver can search freely, and reporting $\tau = [\tilde{\tau}]_{[l,u]}$ at the root $\tilde{\tau}$ recovers an interior optimum. The construction also detects corners: if $f>0$ throughout $[l,u]$ with $\partial f/\partial\tau <0$, the auxiliary problem has its root at $\tilde{\tau} = u+1/a_1$, which bounds to $\tau = u$, and symmetrically $f<0$ with $\partial f/\partial \tau>0$ yields $\tilde{\tau} = l-1/a_0$ and $\tau = l$.

\smalltitle{The grid.}
A political first order condition is not known to be well behaved, so we evaluate it on a grid $\mathcal{T} = \left\lbrace \tau^{(1)},...,\tau^{(M)}\right\rbrace$ with $\tau^{(1)} = l$ and $\tau^{(M)} = u$ rather than trust a solver's path through it. The package can add two outer nodes a hair $\delta>0$ beyond the boundary solutions of \eqref{\refeq:root},
\begin{align}\label{\refeq:extendedGrid}
    \tilde{\tau}^{(0)} = l-\frac{1}{a_0}-\delta, \qquad \tilde{\tau}^{(M+1)} = u+\frac{1}{a_1}+\delta,
\end{align}
so that a pure root search on the grid also returns corners. The solvers of this model do not use them: the selection rule below takes both endpoints of $\mathcal{T}$ as candidates, on the same footing as the crossings.

\smalltitle{Selecting among several roots.}
A zero of $z$ is only a necessary condition for the political optimum: an upward crossing is a local \emph{minimum} of $\mathcal{W}$, several downward crossings may exist, and the optimum may sit at a corner. In every application below $z$ also substitutes a predetermined state consistently with the point of evaluation, the savings shares $s_{t-1,i}/s_{t-1}$ at the tax that savers anticipated. Writing $\mathcal{W}(\tau; D)$ for the objective at a frozen state $D$ and $D(\tau)$ for the consistent one,
\begin{align}\label{\refeq:consistentFOC}
    z(\tau) = \partial_{\tau}\mathcal{W}\left(\tau; D\right)\big|_{D = D(\tau)}.
\end{align}
A zero of $z$ is a stationary point of $\mathcal{W}(\cdot\,; D(\tau^{\star}))$, the object wanted, so roots may be located on the interpolant of $z$. But $z$ is the derivative of no single function of $\tau$: its slope at a crossing, $z' = \mathcal{W}_{\tau\tau} + \mathcal{W}_{\tau D}D'$, is not the second-order condition, and its integral, $\int z\,d\tau = \Delta\left[\mathcal{W}(\tau;D(\tau))\right] - \int \mathcal{W}_D D'\,d\tau$, is neither the objective at a frozen state nor the consistent one. Candidates are therefore tested and ranked at their own states.

\emph{Candidates.} With $\hat{z}$ the piecewise-linear interpolant of $z$ through the nodes of $\mathcal{T}$, the candidate set is
\begin{align}\label{\refeq:candidates}
    \mathcal{C} \equiv \left\lbrace l,u \right\rbrace\cup\left\lbrace \tau\in(l,u)\mbox{ }:\mbox{ }\hat{z}(\tau) = 0\right\rbrace,
\end{align}
every crossing in either direction and the two endpoints; an exact zero at a node, or a run of them between opposite signs, counts as one crossing. Each candidate $c$ comes with its own predetermined state $D_c \equiv D(\tau_c)$, the one savers would have chosen had they anticipated it.

\emph{The objective at a frozen state.} For each candidate the first order condition is re-evaluated on the whole grid with the state frozen at $D_c$, and integrated by the trapezoidal rule,
\begin{align}\label{\refeq:objectiveProfile}
    \mathcal{V}_c(\tau) \equiv \mathcal{W}^{\mathrm{old}}\left(l; D_c\right) + \int_l^{\tau}\partial_{x}\mathcal{W}\left(x; D_c\right)dx, \qquad \tau\in\mathcal{T},
\end{align}
where $\mathcal{W}^{\mathrm{old}}$ is the formal retirees' term, a closed form. $\mathcal{V}_c(\tau)$ equals $\mathcal{W}(\tau; D_c)$ up to terms that do not involve the state and are the same for every candidate, so values compare across candidates as the objectives do. Integrating the condition rather than reading utility levels keeps the smoothed numerical derivatives the condition carries (\cref{\refsec:PEE}); raw levels would bring back the kinks the smoothing removed. Only the formal retirees' term is recomputed, a closed form in $(\tau, D_c)$ given hours, with $D_c$ formed from the candidate's tax and the hours interpolated linearly there.

\emph{The equilibrium test.} A candidate is an \emph{equilibrium} at the state if it is the global maximiser of its own frozen objective,
\begin{align}\label{\refeq:equilibriumTest}
    \mathcal{V}_c(\tau_c) \geq \max_{\tau\in\mathcal{T}}\mathcal{V}_c(\tau),
\end{align}
up to a relative tolerance of $10^{-9}$: the electorate, taking the state as given, prefers $\tau_c$ to every other tax, and the state is the one consistent with $\tau_c$. At a corner $\mathcal{V}_c(\tau_c)$ is the node value. At an interior candidate it is read from the quadratic through the three nearest nodes, and since the crossing is located on $\hat{z}$ and the peak of $\mathcal{V}_c$ on its node values, which differ at second order, two forms at the grid's resolution supplement the reading: the candidate also passes if its nearest node attains the maximum, and, where its stencil touches an infeasible node, if the maximising node lies within one cell of it. A corner's value is read exactly, so neither form applies to it. Passing candidates within one cell of each other are one equilibrium, represented by the interior crossing, which moves continuously with the data. The test decides for every kind of candidate: an upward crossing of $z$ can pass, since its slope is not the second-order condition, and a corner passes exactly when the objective at the corner's own state is maximised there.

\emph{Selection, and what it assumes.} The equilibria at a state need not be unique, and the multiplicity is intrinsic rather than numerical: even when $\mathcal{W}(\cdot\,;D)$ is strictly concave for every $D$, the map $\tau\mapsto\arg\max\mathcal{W}(\cdot\,;D(\tau))$ can have several fixed points, because the shares savers choose depend on the tax they anticipate. When more than one candidate passes we select the one with the highest value of the political objective in its own equilibrium, $\max_c \mathcal{V}_c(\tau_c)$. The selection is part of the definition of the equilibrium computed: the continuation policy that savers and the electorate take as given is built from the selected policy, so the solution is unique only given the rule. The solvers report, for every state and period, the number of candidates, the number of equilibria and whether the fallback below was used, since a selection that switched between equilibria across neighbouring states would put a jump into the policy function the calibration differentiates through. A state where no candidate passes, which the grid can produce at a tangential crossing, is flagged and given the maximiser of $\int_l^{\tau}\hat{z}$ among the endpoints and the downward crossings, so that a recursion is not stopped by one state. When a feasibility mask excludes part of $\mathcal{T}$ (\cref{\refsec:PEE}), all of the above applies to the feasible sub-grid. The state frozen per candidate is the formal retirees' savings shares; under CRRA the savings level is a grid state, constant along the tax grid.

\smalltitle{Root problems that are not maximisations.}
The state fixed points of \cref{\refsec:PEE} are self-consistency conditions in which \emph{any} crossing is a solution, and neither the equilibrium test nor the selection applies to them: the safeguards appropriate to one kind of problem are errors in the other.

\smalltitle{Two implementation conventions.}
Convergence of a gradient-polished solution is judged on $\max\left|z\right|$ against an explicit tolerance, not on the solver's success flag, which standard routines do not report consistently. And a residual measures convergence only at an \emph{interior} optimum: at a corner $z\neq 0$ is the correct answer, so residual diagnostics are restricted to periods and states whose selected policy is interior.
